% =====================================================================
% Chapter 10 -- Conclusion and Future Work
% =====================================================================
\documentclass[../preamble.tex]{subfiles}
\begin{document}

\chapter{Conclusion and Future Work}
\label{ch:conclusion}

\section{Recap of contributions}
\label{sec:conclusion-recap}

This thesis built and evaluated \PEARL{}, a trust-oriented deep-learning
pipeline for PCOS detection from ovarian ultrasound images. The work
contributes:

\begin{enumerate}
  \item a reproducible pipeline that separates deduplication, split
        construction, preprocessing, training, fine-tuning, calibration,
        uncertainty estimation, and explanation;
  \item a controlled SRAD-versus-Gaussian comparison on nine backbones;
  \item a documented letterbox-padding ablation that quantifies a
        cross-dataset shortcut;
  \item a probability-averaged three-architecture ensemble that combines
        DenseNet-121, ConvNeXt-Tiny, and ViT-B/16 for a final external AUC
        of 0.9487 (95\% CI 0.9368--0.9596) on the held-out PCOSgen test set;
  \item two-pass temperature scaling that reduces ensemble Expected
        Calibration Error from 0.0854 to 0.0810;
  \item MC-Dropout predictive-entropy summaries that expose confident and
        uncertain predictions separately; and
  \item Grad-CAM inspection that links the model's attention to anatomical
        regions and to failure modes.
\end{enumerate}

The thesis also documents the difference between the v3 pipeline
specification and the executed final run, so future work can reproduce or
extend the missing components.

\section{Future work}
\label{sec:conclusion-future}

\begin{itemize}
  \item \textbf{Prospective validation.} The most important extension is a
        prospective evaluation at one or more clinical sites with verified
        patient-level grouping, acquisition metadata, and histopathological
        ground truth. The current 0.9487 external AUC is informative but does
        not establish a clinical performance claim.
  \item \textbf{Larger HPO cohorts.} A larger or nested validation set, or
        repeated cross-validation, would stabilise the HPO step. The
        pre-HPO recommendation in this thesis reflects the small validation
        set rather than a permanent choice.
  \item \textbf{Multi-modal fusion.} The pipeline currently uses images only.
        Combining the ultrasound with hormonal and clinical variables would
        move closer to a clinically useful diagnostic workflow, as suggested
        by the cross-attention approach of Lakshmi and Pushpa
        \cite{lakshmi2026}.
  \item \textbf{Federated learning.} A federated training scheme across
        hospitals could address privacy concerns that limit centralised data
        collection. Such a scheme would also enable a much larger
        evaluation cohort.
  \item \textbf{Richer explainability.} LRP, SHAP, and objective
        perturbation-based evaluation were specified in the methodology
        document but not executed \cite{samek2017,patil2025}. Adding them
        and running reader studies would establish whether Grad-CAM, LRP,
        and SHAP produce explanations that clinicians actually trust.
  \item \textbf{Distribution-shift robustness.} The thesis shows that simple
        preprocessing shortcuts produce large domain gaps. Robustness
        techniques such as test-time adaptation, source-free domain
        adaptation, or learned augmentations could close the gap more
        systematically.
  \item \textbf{Class conditioning on the deployment cohort.} The
        sensitivity-targeted operating point is calibrated on the PCOSgen
        test set. Different deployment populations will need different
        thresholds. A deployment-time threshold routine would let a
        hospital tune the operating point to its own prevalence and referral
        cost structure.
\end{itemize}

\section{Closing remark}
\label{sec:conclusion-closing}

Trust-oriented deep learning for medical imaging is not simply a matter of
reported accuracy. It requires reporting what the model does on a different
cohort, what the model believes about its own confidence, and which image
regions drove the decision. The present pipeline goes some way toward each
of those, but the appendices also make clear how much remains. Future work
should treat those gaps as a research agenda rather than as silent
limitations.

\end{document}